\chapter{Residual Boundary of the Arithmetic Branch}
\label{v4-arith:w10-final:chapter}

A Hilbert space with a scaling action and the Weil pairing on
arithmetic test functions are different data. A unitary comparison can
transport a scaling operator to a proposed boundary space. Positivity
of the full Weil pairing instead constrains the zeros of the completed
zeta function. This chapter separates these two assertions before
listing the additional geometric and arithmetic comparisons.

\section{Hilbert Descent, Scaling, and the Weil Pairing}
\label{v4-arith:w10-final:w10-a}

Write $\mathscr A=C_c^\infty(\mathbb R)$ for the compactly supported smooth
functions of the logarithmic variable, with convolution and
$f^*(t)=\overline{f(-t)}$. The Fourier transform
\[
 (\mathcal F_+f)(z)=\int_{\mathbb R}f(t)e^{izt}\,dt
\]
identifies this algebra with its entire image
$\mathscr P=\mathcal F_+\mathscr A$, equipped with pointwise
multiplication and $h^\sharp(z)=\overline{h(\bar z)}$.
Both spaces and their transported topologies are those of
Definition~\ref{v4-arith:w5-a:def:PW}. On the compact-test algebra the signed explicit
formula is
\[
 W(f)=P(f)-W_\infty(f)-W_{\mathrm{fin}}(f),
\]
with the polar, archimedean, and prime terms defined in
Theorem~\ref{v4-arith:w5-a:thm:weil-explicit}. The Hermitian form is
$q_W(f,g)=W(f*g^*)$. None of these definitions uses a boundary
Hilbert-space construction.

For comparison, the continuous scaling model is
\[
 H_c=L^2(\mathbb R_{>0},dx/x),\qquad
 (U_a\phi)(x)=\phi(e^{-a}x).
\]
Under $t=\log x$, put $D_c=-i\partial_t$ on
$H^1(\mathbb R)$ and $\Theta_c=\tfrac12+iD_c$ on the same domain.
Here $H^1$ consists of square-integrable functions whose weak first
derivative is square-integrable. The translation group satisfies
$U_a=e^{-iaD_c}$, and
\[
 e^{-a\Theta_c}=e^{-a/2}U_a,\qquad
 |e^{-a\Theta_c}|=e^{-a/2}1_{H_c}\quad(a\in\mathbb R).
\]
This norm identity is unconditional and makes no assertion about
$q_W$.

\begin{theorem}[arithmetic positivity and boundary transport]
\label{v4-arith:w10-final:thm:cs-descent-atp-rh}
The archimedean Tate positivity (A-TP) statement of
Conjecture~\ref{v4-arith:w5-a:conj:archimedean-tate-positivity} is equivalent
to RH. No Hilbert-descent hypothesis is needed for this equivalence.

Suppose separately that a Hilbert-descent construction supplies a
space $H_B$. Specify a closed normal operator $\Theta_B$ and a unitary
$V_\Sigma:H_B\to H_c$ such that
\[
 V_\Sigma\operatorname{Dom}(\Theta_B)=\operatorname{Dom}(\Theta_c),
 \qquad V_\Sigma\Theta_B=\Theta_cV_\Sigma
 \quad\text{on }\operatorname{Dom}(\Theta_B).
\]
Then every $e^{-a\Theta_B}$, $a\in\mathbb R$, is bounded and
\[
 e^{-a\Theta_B}=V_\Sigma^{-1}e^{-a\Theta_c}V_\Sigma,
 \qquad |e^{-a\Theta_B}|=e^{-a/2}1_{H_B}.
\]
Thus descent together with A-TP implies RH through A-TP; the surface
operator comparison additionally requires the stated unitary and
generator-domain data.
\end{theorem}

\begin{proof}
The signed explicit formula gives
\[
 q_W(f,f)=P(f*f^*)-W_\infty(f*f^*)-W_{\mathrm{fin}}(f*f^*).
\]
Consequently A-TP is precisely $q_W(f,f)\ge0$ for every
$f\in\mathscr A$. Apply
Corollary~\ref{v4-arith:w5-a:cor:weil-positivity-RH}.
For the operator statement, the domain and intertwining equalities
identify the spectral resolutions of $\Theta_B$ and $\Theta_c$.
The latter has spectrum in $1/2+i\mathbb R$, on which
$z\mapsto e^{-az}$ has constant modulus $e^{-a/2}$.
Normal functional calculus proves both identities.
This is the comparison specified in
Theorem~\ref{v4-arith:w10-a:thm:wilson-surface-compat}.
\end{proof}

\begin{remark}[one exponential does not determine its generator]
\label{v4-arith:w10-final:rem:w10-a-shape}
The Hilbert limit in
Conjecture~\ref{v4-arith:w10-a:conj:condensed-descent} and the comparison
$V_\Sigma$ are separate inputs. Even equality of the surface operator
at $a=1$ does not replace the generator comparison:
$\Theta_c$ and $\Theta_c+2\pi i\,1$ have the same domain and the same
exponential at $a=1$, but are distinct operators. They give distinct
one-parameter families at nonintegral $a$. To identify a native
surface theory one must in addition relate its independently defined
observable to this transported operator. No map from arithmetic test
functions or arithmetic cohomology is supplied by the norm identity.
\end{remark}

The radial spherical transform of Theorem~\ref{v4-arith:radial:transform}
provides a different unitary comparison: tree adjacency becomes
multiplication by $2\sqrt p\cos\theta$, and the root state becomes
integration against $\nu_p$. Section~\ref{v4-arith:radial:finite}
extends this comparison to finite sets of primes. An arithmetic
realization still requires a map from the primitive trace carrier
preserving the Hecke operations and pairing. Tame residue and Gersten
operations require additional maps on common coefficients. The positive
tree state supplies no identification with the full Weil functional
$W$ defined above.

\section{The Full Weil Positivity Condition}
\label{v4-arith:w10-final:tier1}

In compact coordinates the arithmetic condition is
\begin{equation}
\label{v4-arith:w10-final:eq:atp}
 W(f*f^*)\ge0\qquad(f\in\mathscr A).
\end{equation}
Define $\mathscr W:\mathscr P\to\mathbb C$ by
$\mathscr W(h)=W(\mathcal F_+^{-1}h)$. Its equivalent Fourier form is
\[
 \mathscr W(hh^\sharp)\ge0\qquad(h\in\mathscr P).
\]
The product here is pointwise multiplication of entire functions,
not convolution on $\mathscr P$.

\begin{proposition}[the two test-space forms of Weil positivity]
\label{v4-arith:w10-final:prop:tier1}
The compact and Fourier inequalities above are equivalent, and each
is equivalent to RH. Restricting to an arbitrary proper test subspace
requires a separate argument before the same criterion can be used.
\end{proposition}

\begin{proof}
Fubini gives
$\mathcal F_+(f*f^*)=(\mathcal F_+f)(\mathcal F_+f)^\sharp$.
By definition $\mathcal F_+:\mathscr A\to\mathscr P$ is onto;
Fourier inversion makes it one-to-one. Hence the two universally
quantified inequalities agree. Their equivalence with RH is
Corollary~\ref{v4-arith:w5-a:cor:weil-positivity-RH} on exactly this test
space. Neither an auxiliary digamma quadratic form nor continuous
scaling purity replaces the full polar-minus-archimedean-minus-prime
expression.
\end{proof}

\section{Comparison and Descent Conjectures}
\label{v4-arith:w10-final:tier2}

\begin{definition}[comparison and descent inputs]
\label{v4-arith:w10-final:def:prior-table}
\ClaimStatusDefinitional The comparison and descent boundary contains
the following named inputs.
\begin{center}
\small
\begin{tabular}{p{3.6cm}p{7.8cm}p{1.4cm}}
\toprule
\textbf{Input} & \textbf{Carries} & \textbf{Form}\\
\midrule
Wild-EG enhancement &
F1.a.2.class-match;
\(\mathrm{L}_{\mathrm{ACD}}\) Eisenstein \(\mathrm{Ext}^{1}\) at
\(\{0,1\}\); BC-glue wild inertia &
\(\mathrm{open}\)\\
Cross-prime transition and effective descent &
BC-diamond-glue cross-prime prismatic glue &
\(\mathrm{open}\)\\
Arithmetic AG singular support &
global arithmetic Arinkin--Gaitsgory singular-support condition and
the centre of BC-diamond-glue &
\(\mathrm{open}\)\\
Clausen--Scholze condensed-Hilbert descent &
Quantum arithmetic Chern--Simons bulk extension &
\(\mathrm{open}\)\\
Consani--Marcolli adele-class duality &
Positselski \(G\)-row duality and the abelian Heisenberg base-change
line &
\(\mathrm{open}\)\\
Chenevier Selmer-ring, trace and equivariant-descent datum &
Pairwise sheaf comparison under the full Ch1 hypotheses &
\(\mathrm{open}\)\\
BCglue centre-compatible categorical functors and coefficient/family comparisons &
K5 and the Jacquet--Shalika trace comparison in BC-diamond-glue &
\(\mathrm{open}\)\\
Ben-Zvi lift at infinity &
\(\mathrm{L}_{\mathrm{ACD}}\) archimedean compatibility and finite
uniformity &
\(\mathrm{open}\)\\
BZN wild trace data &
Large-prime semisimplicity and the wild \(\mathrm{Ext}\)-trace pairing &
\(\mathrm{open}\)\\
Symmetric-square weight limit &
Weil--Petersson weight-aspect positivity for holomorphic cusp forms &
\(\mathrm{open}\)\\
Arithmetic chiral bar nilpotency &
The square-zero property of the proposed arithmetic chiral
pre-differential \(d^{\mathrm{Ar}}\) &
\(\mathrm{open}\)\\
\bottomrule
\end{tabular}
\end{center}
\end{definition}

The relative Galois bar in Chapter~\ref{v4-arith:relative-bar:chapter}
has a square-zero differential over each $\mathbb Z/\ell^r$.
Theorem~\ref{v4-arith:endpoint:thm:main} compares its endpoint products
in every arity, and Corollary~\ref{v4-arith:endpoint:cor:adic} passes
to the specified restricted adic limit. The last row concerns the
separate arithmetic chiral pre-differential. Identifying its carrier,
residue operations, and completion with the relative Galois construction
remains an additional comparison problem.

\begin{theorem}[visible non-RH residuals]
\label{v4-arith:w10-final:thm:four-prior}
\ClaimStatusProvedHere Each input in
Definition~\ref{v4-arith:w10-final:def:prior-table} is a non-RH residual named
elsewhere in the arithmetic branch. The statement does not assert
exhaustion of all possible residuals.
\end{theorem}

\begin{proof}
With its stated comparison hypotheses, Wild-EG (the Emerton--Gee
wild-inertia enhancement) is an input to the residual-characteristic
local comparison for F1.a.2 in Theorem~\ref{v4-arith:w10-b:thm:class-match}.
The global exceptional Eisenstein comparison in
Remark~\ref{v4-arith:w10-d:thm:class-match} additionally requires a common
coefficient realization, the automorphic carriers, a residue-to-Yoneda
comparison, and compatible full derived maps.
For BC-diamond-glue it must be accompanied
by the block-centre maps, compatible
functors and normalized class comparisons
specified in
Proposition~\ref{v4-arith:w10-e:prop:ell-equals-v-wild}.
The ordinary \(\ell\neq p\) scalar comparison
does not supply these data.

The cross-prime input is the specified transition, coherent-support,
filtration and higher-coherence datum of
Hypothesis~\ref{v4-arith:bcglue-k5:lem:prismatic-flatness}, together
with the effective descent required by its gluing theorem. It is
not an independently proved pro-\'etale or prismatic-site flatness
statement. Arithmetic AG (arithmetic
Arinkin--Gaitsgory singular support) is the global singular-support
conjecture~\ref{v4-arith:conj:F1-global-AG}, required at the centre
of BC-diamond-glue by Corollary~\ref{v4-arith:w10-e:cor:bcglue-final-updated}.
The specified condensed-Hilbert descent datum supplies the Hilbert,
trace, and finite-place Wilson clauses of the chosen quantum ACS model
by Theorem~\ref{v4-arith:w10-a:thm:descent-equivalent}. The archimedean
operator comparison also requires the unitary and generator domains in
Theorem~\ref{v4-arith:w10-a:thm:wilson-surface-compat}.
The other rows are the named conditional inputs constructed in
Chapters~\ref{v4-arith:w5-d:chapter},
\ref{v4-arith:w5-h:bcglue-core-and-F2c},
\ref{v4-arith:bcglue-k5},
\ref{v4-arith:lacd-sub}, \ref{v4-arith:bzn-rcompletion},
\ref{v4-arith:w6-b:chapter}, and
\ref{v4-arith:arakelov-bar}.  These are stated as separate comparison hypotheses. No implication
from A-TP to these constructions, or in the reverse direction, is
established by the Weil positivity criterion.
\end{proof}

\begin{remark}[Wild-EG concentration]
\label{v4-arith:w10-final:rem:wild-eg}
Wild-EG is a common input to the three comparison problems
F1.a.2, the exceptional $\mathrm{L}_{\mathrm{ACD}}$ Eisenstein
$\mathrm{Ext}^1$ crossing at $\{0,1\}$, and wild-inertia
BC-diamond-glue. Each retains its stated class-comparison hypotheses.
For BC-diamond-glue, one must also construct the fixed-determinant
block-ring isomorphisms, compatible fully faithful functors and
normalized class maps of
Proposition~\ref{v4-arith:w10-e:prop:ell-equals-v-wild}, together with the
coefficient, family and global comparisons. The wild-EG enhancement
alone does not construct these data.
\end{remark}

\section{The Generalised Ramanujan Conjecture for
\texorpdfstring{\(\mathrm{GL}_{4}\)}{GL\_4}}
\label{v4-arith:w10-final:sub-rh}

The Rankin--Selberg square of a Maass form on \(\mathrm{GL}_{2}/\mathbb{Q}\)
comes from an isobaric automorphic object on
\(\mathrm{GL}_{4}/\mathbb{Q}\). Full temperedness of that
\(\mathrm{GL}_{4}\) object at every finite place, equivalently GRC for
its cuspidal factors, closes the Maass sector of the
\(\mathcal{V}^{\mathrm{prim}}\) Hochschild-trace identity. This
conjecture is separated from A-TP in the proof: it concerns Satake
eigenvalues, not zero placement.

\begin{hypothesis}[GRC for \(\mathrm{GL}_{4}/\mathbb{Q}\) in the Selberg class]
\label{v4-arith:w10-final:hyp:GRC-GL4}
\ClaimStatusConditional The isobaric automorphic representations of
\(\mathrm{GL}_{4}/\mathbb{Q}\) arising from the Rankin--Selberg square
of a Maass form have tempered cuspidal factors at every finite place.
\end{hypothesis}

\begin{proposition}[GRC-\(\mathrm{GL}_{4}\) closes the Maass sector]
\label{v4-arith:w10-final:prop:GRC-GL4}
\ClaimStatusProvedHereConditional Conditional on
Hypothesis~\ref{v4-arith:w10-final:hyp:GRC-GL4}, full Ramanujan holds for the
Maass forms entering the \(\mathcal{V}^{\mathrm{prim}}\)
Hochschild-trace \(P_{3}\) identity of Theorem~\ref{v4-arith:w10-c:thm:main}.
This closes the Maass sector. No equivalence with RH is asserted here; no implication between this
hypothesis and A-TP is used.
\end{proposition}

\begin{proof}
Chapter~\ref{v4-arith:w10-c:chapter} proves that the
Arthur--Ramakrishnan--GRC hypothesis implies full Ramanujan for Maass forms
(Theorem~\ref{v4-arith:w10-c:thm:AR-implies-Ramanujan}) and reduces that
hypothesis to Hypothesis~\ref{v4-arith:w10-final:hyp:GRC-GL4}
(Proposition~\ref{v4-arith:w10-c:prop:sub-RH-strict}). The \(P_{3}\) identity
is the exact conditional statement of Theorem~\ref{v4-arith:w10-c:thm:main}.
A-TP concerns the zeros of \(\zeta\) through Weil positivity;
GRC-\(\mathrm{GL}_{4}\) concerns Satake parameters.  The argument uses
the two residuals as separate inputs.
\end{proof}

\section{\(F_{1}\text{.a.1}\) Small Residuals}
\label{v4-arith:w10-final:tier3}

Two small residuals occur in the \(F_{1}\text{.a.1}\) table of
Chapter~\ref{v4-arith:w10-f:chapter}. They are not a count of all
small residuals in the arithmetic branch.

\begin{proposition}[small residuals]
\label{v4-arith:w10-final:prop:small-residuals}
\ClaimStatusProvedHere Besides the generator-comparison inputs and the
automorphy/patching hypotheses isolated in
Lemma~\ref{v4-arith:w10-f:lem:BLGG-descent} and
Lemma~\ref{v4-arith:w10-f:lem:Emerton-patching},
the \(F_{1}\text{.a.1}\) table has the following two small residual tails:
\begin{enumerate}
\item \textbf{Supersingular small-weight tail}
(Conjecture~\ref{v4-arith:w10-f:conj:ss-small-weight}): non-Fontaine--Laffaille
representations with Hodge--Tate weights
\(k \in \{p-1, p\}\), requires a non-formal theorem.
\item \textbf{\(A_{5}\) non-solvable descent defect}
(Conjecture~\ref{v4-arith:w10-f:conj:ns-defect}): local-global descent defect
in the icosahedral sector, requires a non-formal theorem.
\end{enumerate}
\end{proposition}

\begin{proof}
The two listed residual tails are read from
Proposition~\ref{v4-arith:w10-f:prop:F1a1-table-w10f}. The Fontaine--Laffaille
supersingular core is reduced to the generator comparison of
Theorem~\ref{v4-arith:w10-f:thm:f1a1-ss-FL}. The odd-determinant icosahedral
core is reduced to the conditional theorem of
Theorem~\ref{v4-arith:w10-f:thm:A5-odd-closure}. The two tails listed above are
additional obstructions in that decomposition.
\end{proof}

\section{Arithmetic Hypotheses and Realization Data}
\label{v4-arith:w10-final:tier4}

\begin{remark}[scope of the arithmetic conditions]
\label{v4-arith:w10-final:prop:hziv-separate}
The conditions above concern arithmetic comparison, descent,
temperedness and small-prime regularity. A marked realization as in
\Cref{v4-ch:marked-realizations} additionally requires its source and
target, coefficient functor, operation-preserving maps and evaluation
identity. The arithmetic conditions alone construct none of those
maps and imply no realization assertion for an object in another
coefficient category.
\end{remark}

\section{Residual Boundary}
\label{v4-arith:w10-final:residual-boundary}

\begin{theorem}[arithmetic branch obstruction decomposition]
\label{v4-arith:w10-final:thm:near-closure}
\ClaimStatusProvedHere The arithmetic branch has, at least, the
following obstructions:
\begin{enumerate}
\item one statement equivalent to RH: archimedean Tate positivity;
\item the comparison and descent inputs of
Definition~\ref{v4-arith:w10-final:def:prior-table};
\item one Ramanujan-type hypothesis separated from A-TP: GRC for
\(\mathrm{GL}_{4}/\mathbb{Q}\);
\item the two \(F_{1}\text{.a.1}\) small residuals of
Proposition~\ref{v4-arith:w10-final:prop:small-residuals}.
\end{enumerate}
The arithmetic conditions listed here do not construct the marked
realizations of Chapter~\ref{v4-ch:marked-realizations}.
\end{theorem}

\begin{proof}
Item (1) is Proposition~\ref{v4-arith:w10-final:prop:tier1}. Item (2) is
Theorem~\ref{v4-arith:w10-final:thm:four-prior}. Item (3) is
Proposition~\ref{v4-arith:w10-final:prop:GRC-GL4}. Item (4) is
Proposition~\ref{v4-arith:w10-final:prop:small-residuals}. A marked realization
requires the separate data described in
Remark~\ref{v4-arith:w10-final:prop:hziv-separate}; the listed arithmetic
conditions do not supply those data.
No implication from the table to RH is asserted except through A-TP.
\end{proof}

\begin{remark}[domain of the boundary]
\label{v4-arith:w10-final:rem:not-rh}
The theorem isolates the displayed RH-equivalent content into A-TP and
places the visible non-RH obstructions in named inputs.  It does not
prove RH and it does not claim that the arithmetic comparison is
complete.
\end{remark}
